\subsection{Analisis Sistem DecMed}
\label{subsec:analisis-sistem-decmed}

Sebagaimana dijelaskan pada Bab \ref{sec:decmed}, DecMed adalah sistem manajemen RME terdesentralisasi dengan memanfaatkan IOTA sebagai penyimpanan \textit{on-chain}, IPFS sebagai penyimpanan \textit{off-chain}, serta PRE untuk mendukung distribusi akses data yang aman. Arsitektur ini sudah memberikan fondasi untuk menjaga integritas data dan mengurangi ketergantungan pada penyimpanan terpusat. Namun, mekanisme kontrol akses yang digunakan masih bertumpu pada token kapabilitas statis, sehingga belum sepenuhnya mendukung kebutuhan otorisasi yang lebih rinci pada skenario layanan kesehatan yang kolaboratif.

Keterbatasan pertama terletak pada representasi hak akses yang masih bersifat \textit{coarse-grained}. Token akses pada DecMed belum dirancang untuk membawa batasan kontekstual, misalnya jenis data yang boleh diakses, aksi yang diizinkan, masa berlaku akses, atau batasan delegasi lanjutan. Padahal, pada praktik pelayanan kesehatan, setiap tenaga medis tidak selalu membutuhkan keseluruhan isi RME pasien. Kondisi ini menunjukkan bahwa mekanisme akses kontrol pada DecMed belum cukup untuk menjawab kebutuhan kontrol akses granular sebagaimana dirumuskan pada tujuan pertama Tugas Akhir ini.

Apabila dibandingkan dengan karakteristik Macaroons pada Subbab \ref{sec:macaroons}, keterbatasan tersebut menjadi semakin jelas. Token pada DecMed pada dasarnya hanya merepresentasikan hak akses yang telah ditetapkan saat token diterbitkan, sehingga sulit menambahkan batasan baru secara aman ketika akses perlu diteruskan kepada pihak lain. Sebaliknya, Macaroons menyediakan struktur \textit{chained-HMAC} dan \textit{caveat} yang memungkinkan setiap proses delegasi sekaligus membawa batasan tambahan, misalnya batasan terhadap jenis data, masa berlaku, atau konteks penggunaan. Oleh karena itu, pemanfaatan Macaroons bukan sekadar pilihan implementasi alternatif, melainkan kebutuhan teknis untuk memenuhi tuntutan otorisasi granular dan delegasi terbatas.

Keterbatasan kedua berkaitan dengan organisasi data pada lapisan \textit{off-chain}. DecMed memang telah memisahkan penyimpanan \textit{on-chain} dan \textit{off-chain}, tetapi pendekatan tersebut belum secara eksplisit menunjukkan segmentasi data klinis yang cukup rinci untuk menegakkan kebijakan akses \textit{fine-grained}. Jika data klinis masih dikelola sebagai objek terenkripsi yang terlalu besar atau terlalu umum, maka pemberian akses kepada seorang tenaga medis berisiko membuka cakupan data yang lebih luas daripada yang dibutuhkan. Dengan demikian, diperlukan rancangan penyimpanan \textit{off-chain} yang tersegmentasi agar granularitas kebijakan dapat ditegakkan hingga level himpunan data yang relevan.

Keterbatasan ketiga muncul pada proses pendelegasian akses. Pada DecMed, pasien menjadi pihak yang harus terlibat saat akses baru perlu diterbitkan untuk tenaga medis lain. Asumsi ini kurang sesuai dengan alur pelayanan kesehatan yang bersifat dinamis, dimana dokter penanggung jawab pasien dapat perlu mendelegasikan sebagian hak akses kepada residen, perawat, atau unit laboratorium tanpa harus menunggu pasien kembali \textit{online}. Oleh sebab itu, DecMed memerlukan mekanisme delegasi berantai yang tetap aman tetapi dapat dilakukan secara luring melalui proses atenuasi hak akses. Kebutuhan inilah yang menjadi dasar pemanfaatan Macaroons pada penelitian ini.

Kelemahan fungsional pada granularitas, segmentasi data, dan pendelegasian ini tidak hanya berdampak pada efisiensi operasional medis, tetapi juga membuka celah keamanan yang signifikan. Untuk memetakan risiko tersebut secara sistematis, dilakukan analisis ancaman menggunakan kerangka kerja STRIDE pada subbab berikutnya.
